\chapter{关键公式与约定}

本附录给出主文反复使用的公式索引。页码不进入公式表，完整来源与读取深度见B04
批次的`coverage.tsv`和`notes.md`。

\begin{longtable}{@{}p{0.18\textwidth}p{0.40\textwidth}p{0.34\textwidth}@{}}
  \caption{B04关键公式与必须同时声明的约定}\label{tab:b04-formula-reference}\\
  \toprule
  对象 & 公式或定义 & 配套约定与限制 \\
  \midrule
  \endfirsthead
  \toprule
  对象 & 公式或定义 & 配套约定与限制 \\
  \midrule
  \endhead
  \bottomrule
  \endfoot
  光锥坐标 & $v^\pm=(v^0\pm v^3)/\sqrt2$，$g_{+-}=1$ & 与$P^\pm$、Fourier相位和Wilson线方向一起固定 \\
  夸克PDF & 式\eqref{eq:b04-lc-quark} & 指定场序、$+$规范链、$\overline{\rm MS}$或格点方案 \\
  胶子PDF & 式\eqref{eq:b04-lc-gluon} & 左侧为$xg$；$F_{+}^{\ \mu}$指标和伴随链接须一致 \\
  守恒 & 式\eqref{eq:b04-pdf-sumrules} & 动量规则含胶子；价数规则按味和反夸克定义 \\
  DGLAP & 式\eqref{eq:b04-dglap-nonsinglet}与\eqref{eq:b04-dglap-singlet} & $\mu$为因子化标度；端点采用$[\,]_+$减除 \\
  劈裂核 & 式\eqref{eq:b04-splitting} & LO；$n_f$为活跃轻味数，颜色因子按$SU(3)$ \\
  Mellin & 式\eqref{eq:b04-mellin}至\eqref{eq:b04-gamma-qq} & 矩指标$n$与算符自旋、多项式次数同步 \\
  OPE & 式\eqref{eq:b04-short-ope}与\eqref{eq:b04-lightray-ope} & 普通短距与light-ray的幂次计数不同 \\
  准PDF & 式\eqref{eq:b04-quasi}与\eqref{eq:b04-lamet-factorization} & 必须声明$\Gamma$、投影、$P_z$和匹配方案 \\
  ITD & 式\eqref{eq:b04-itd}与\eqref{eq:b04-ritd} & $\nu=zP_z$；分母为零动量/真空同几何矩阵元 \\
  胶子算符 & 式\eqref{eq:b04-gluon-combination} & 只适用于所列指标组合；不可外推到全部36分量 \\
  GPD & 式\eqref{eq:b04-gpd-kinematics}与\eqref{eq:b04-vector-gpd} & $\Delta=P_f-P_i$；旋量、相位和$E$项整体固定 \\
  GPD求和 & 式\eqref{eq:b04-gpd-formfactors}与\eqref{eq:b04-ji-sumrule} & leading twist；形状因子与$t\to0$外推不可省略 \\
\end{longtable}

\chapter{计算验收清单}

本清单把“代码可运行”和“物理量可发布”分开。建议每个结果至少标注
\texttt{diagnostic-only}、\texttt{candidate-physics}或\texttt{accepted}，并保存
对应证据。

\begin{enumerate}
  \item \textbf{约定：}固定光锥坐标、度规、Fourier相位、场强符号、生成元归一化、
  Euclidean gamma约定和反夸克定义。
  \item \textbf{算符：}声明Dirac结构、味组合、色表示、Wilson线端点与方向、
  胶子指标缩并和投影。$\gamma^t$与$\gamma^z$不得在未处理连续混合时互换。
  \item \textbf{归一化：}验证$x=0$、$z=0$、零动量和真空减除；检查动量求和、
  价数、PDF前向极限和GPD形状因子极限。
  \item \textbf{UV：}证明线性Wilson线发散被同几何减除，剩余对数进入正确方案；
  对混合通道给出完整矩阵，而不是一个标量$Z$。
  \item \textbf{匹配：}固定匹配阶数、$\mu$、投影、方案转换、味单态矩阵和
  $[\,]_+$端点。扫描$\mu$与阶数，检查残差是否进入误差预算。
  \item \textbf{幂修正：}先声明目标是twist-2、twist-3还是更高；分别报告端点增强、
  靶质量和重正子模糊。不得把式\eqref{eq:b04-qpdf-power}的$1/P_z^2$自动解释为
  twist-4。
  \item \textbf{Fourier与ITD：}报告$z_{\max}$、$\lambda=zP_z$范围、尾部模型、
  相位、归一化和端点权。大$z$外推结果必须与不同窗口和不同$P_z$交叉检验。
  \item \textbf{极限：}连续极限、无穷动量、手征、体积和激发态至少各有两个可辨识
  参数点；单点结果只能作为该点的有效结果。
  \item \textbf{误差：}传播统计协方差和系统误差相关性；$z_s$、$z_0$、$\mu$、
  $\lambda$窗口和匹配阶数都应作为独立系统源。
  \item \textbf{证据：}保留组态标识、作用量与流参数、源型、缓存键、真残差、
  原始或可重建矩阵元及哈希。中心值一致不能替代窗口和方案检查。
\end{enumerate}

\paragraph{物理接受条件.}
只有当算符与约定闭合、UV/rapidity及混合受控、匹配和幂修正处于声明精度、
所有极限在窗口内稳定、误差包含相关系统项，并与独立方案或独立数据交叉一致时，
才可把结果标记为\texttt{accepted}。若仅完成矩阵元、比值或单参数点，应标记为
\texttt{diagnostic-only}；若主要窗口已通过但连续极限或单态混合尚未闭合，应标记
为\texttt{candidate-physics}，不得改称最终物理PDF或GPD。
